%% abstrak-en.tex
%% Abstrak dalam Bahasa Inggris (Abstract)
%% Format sesuai Panduan Penyusunan Tesis ITS
%%
%% EDIT DATA DI BAWAH INI SESUAI DATA TESIS ANDA

\newpage
\thispagestyle{plain}

%% ===== JUDUL ABSTRAK =====
\begin{center}
    {\fontsize{14}{16}\selectfont\bfseries JUDUL ABSTRAK BAHASA INGGRIS}
\end{center}

\vspace{0.8cm}

%% ===== DATA MAHASISWA =====
\noindent
\begin{tabular}{@{}p{4.5cm}@{}p{0.3cm}@{}p{9cm}@{}}
    Name & : & \\
    Student Identity Number & : & \\
    Supervisor & : & \\
\end{tabular}

\vspace{1cm}

%% ===== HEADING ABSTRAK =====
\begin{center}
    {\fontsize{12}{14}\selectfont\bfseries ABSTRACT}
\end{center}

\vspace{0.5cm}

%% ===== ISI ABSTRAK =====
%% Ganti teks di bawah ini dengan isi abstrak Anda
%% Hapus \colorbox{yellow}{} dan \sout{} setelah mengisi data sebenarnya
%% Maksimum 350 kata

\noindent
\hspace{1.5cm}Lorem ipsum

\hspace{1.5cm}Maksimum 350 kata.

\vspace{0.5cm}

%% ===== KATA KUNCI =====
\noindent
\textbf{Keywords:} ditulis urut abjad